% -------------------------------------------------------------------------------------
\section{模式与幂次计数}\label{sec:scet-modes}

有效场论的第一步是决定保留哪些自由度。对一个由硬标度 $Q$ 与小参数 $\lambda$ 刻画的观测量，第~4~章的钳夹分析给出领头区域；SCET 为每个领头区域引入一个场。

\begin{definition}[模式]\label{def:scet-modes}
设 $Q$ 为硬标度，$\lambda\ll1$ 为由观测量确定的幂次计数参数（推力：$\lambda=\sqrt\tau$；横动量分布：$\lambda=q_T/Q$）。下表中的动量区域称为 SCET 的\emph{模式}（mode）：
\begin{center}\small
\renewcommand{\arraystretch}{1.15}
\begin{tabular}{@{}lccl@{}}
\toprule
模式 & $(\bar n\cdot p,\ n\cdot p,\ p_\perp)/Q$ & $p^2$ & 场 \\
\midrule
硬 & $(1,1,1)$ & $Q^2$ & 积出为匹配系数 \\
$n$ 共线 & $(1,\lambda^2,\lambda)$ & $\lambda^2Q^2$ & $\xi_n$，$A_n$ \\
$\bar n$ 共线 & $(\lambda^2,1,\lambda)$ & $\lambda^2Q^2$ & $\xi_{\bar n}$，$A_{\bar n}$ \\
超软 & $(\lambda^2,\lambda^2,\lambda^2)$ & $\lambda^4Q^2$ & $q_{us}$，$A_{us}$ \\
软 & $(\lambda,\lambda,\lambda)$ & $\lambda^2Q^2$ & $q_s$，$A_s$ \\
硬共线 & $(1,\lambda,\lambda^{1/2})$ & $\lambda Q^2$ & 在 SCET$_{\rm I}\to$SCET$_{\rm II}$ 中积出 \\
Glauber $n\bar n$ & $(\lambda^2,\lambda^2,\lambda)$ & $\simeq-\bm p_{\mathrm T}^2$ & 离壳，不是传播的模式 \\
Glauber $ns$ & $(\lambda,\lambda^2,\lambda)$ & $\simeq-\bm p_{\mathrm T}^2$ & 同上 \\
Glauber $\bar ns$ & $(\lambda^2,\lambda,\lambda)$ & $\simeq-\bm p_{\mathrm T}^2$ & 同上 \\
\bottomrule
\end{tabular}
\end{center}
% VF-SCET-1
只含硬、两个共线与超软模式的理论称为 SCET$_{\rm I}$；只含硬、两个共线（虚度 $\lambda^2Q^2$）与软模式的理论称为 SCET$_{\rm II}$。由 $(\bar n\cdot p,n\cdot p)=\sqrt2(p^+,p^-)$，表中的标度与\cref{def:ir-scalings}中 $(k^+,k^-,k_T)$ 的标度逐项相同；Libby--Sterman 意义下的“软”即此处的超软。
\uses{def:ir-scalings}
\end{definition}

Rothstein--Stewart 把 Glauber 动量写作 $Q(\lambda^a,\lambda^b,\lambda)$（按 $(n\cdot p,\bar n\cdot p,p_\perp)$ 排列），$a+b>2$，$\{a,b\}=\{2,2\},\{2,1\},\{1,2\}$ 分别称为 $n$–$\bar n$、软–$n$、软–$\bar n$ Glauber\cite[表~1]{RothsteinStewart2016}；交换前两项后即上表的后三行。\Cref{fig:scet-modes}把这些模式画在 $(\bar n\cdot p,n\cdot p)$ 的对数平面上：在壳模式落在等虚度线上，Glauber 模式的光锥分量与某个在壳模式相同，但横动量大得多，因而远离质壳。

%%% BEGIN fig:scet-modes
\begin{figure}[htbp]
\centering
\begin{tikzpicture}[fd,scale=1.05]
  % axes: a = log_{1/lambda}(Q/nbar.p) to the right, b = log_{1/lambda}(Q/n.p) upwards
  \draw[lcaxis] (-0.25,0) -- (4.75,0);
  \draw[lcaxis] (0,-0.25) -- (0,4.6);
  \node[below,font=\footnotesize] at (3.9,-0.3) {$\log_{1/\lambda}(Q/\bar n\cdot p)$};
  \node[left,font=\footnotesize,align=right] at (-0.08,4.45) {$\log_{1/\lambda}(Q/n\cdot p)$};
  \foreach \t in {1,2,3,4} {\draw (\t,0) -- (\t,-0.07) node[below,font=\scriptsize]{$\t$};
                           \draw (0,\t) -- (-0.07,\t) node[left,font=\scriptsize]{$\t$};}
  % lines of constant virtuality a+b = 1, 2, 4
  \draw[gray!70,dashed,line width=0.5pt] (1.25,-0.25) -- (-0.25,1.25);
  \draw[gray!70,dashed,line width=0.5pt] (2.3,-0.3) -- (-0.3,2.3);
  \draw[gray!70,dashed,line width=0.5pt] (4.3,-0.3) -- (-0.3,4.3);
  \node[gray!50!black,font=\scriptsize,anchor=west] at (4.3,-0.3) {$p^2\sim\lambda^4Q^2$};
  \node[gray!50!black,font=\scriptsize,anchor=west] at (2.3,-0.3) {$\lambda^2Q^2$};
  \node[gray!50!black,font=\scriptsize,anchor=north] at (1.25,-0.32) {$\lambda Q^2$};
  % zero rapidity
  \draw[gray!60,dotted,line width=0.6pt] (0,0) -- (4.2,4.2);
  \node[gray!45!black,font=\scriptsize,rotate=45,anchor=south] at (3.65,3.65) {$y=0$};
  % on-shell modes
  \fill (0,0) circle[radius=2.2pt]; \node[font=\footnotesize,anchor=north east] at (-0.04,-0.04) {硬};
  \fill[blue!70!black] (0,2) circle[radius=2.2pt]; \node[font=\footnotesize,anchor=west] at (0.1,2.08) {$n$ 共线};
  \fill[orange!80!black] (2,0) circle[radius=2.2pt]; \node[font=\footnotesize,anchor=south west] at (2.05,0.05) {$\bar n$ 共线};
  \fill[green!50!black] (1,1) circle[radius=2.2pt]; \node[font=\footnotesize,anchor=south east] at (0.98,1.04) {软};
  \fill[green!50!black] (2,2) circle[radius=2.2pt]; \node[font=\footnotesize,anchor=south west] at (2.12,2.12) {超软};
  \fill[gray] (0,1) circle[radius=1.7pt]; \node[font=\scriptsize,anchor=west] at (0.08,0.93) {硬共线};
  \fill[gray] (1,0) circle[radius=1.7pt];
  % Glauber (off shell): open symbols
  \draw[red!60!black,line width=0.7pt] (2,2) circle[radius=4.6pt];
  \draw[red!60!black,line width=0.7pt,fill=white] (1,2) ++(-3.4pt,0) -- ++(3.4pt,3.4pt) -- ++(3.4pt,-3.4pt) -- ++(-3.4pt,-3.4pt) -- cycle;
  \draw[red!60!black,line width=0.7pt,fill=white] (2,1) ++(-3.4pt,0) -- ++(3.4pt,3.4pt) -- ++(3.4pt,-3.4pt) -- ++(-3.4pt,-3.4pt) -- cycle;
  \node[red!60!black,font=\scriptsize,anchor=south] at (1,2.12) {$G_{ns}$};
  \node[red!60!black,font=\scriptsize,anchor=west] at (2.13,0.98) {$G_{\bar ns}$};
  \node[red!60!black,font=\scriptsize,anchor=north west] at (2.12,1.9) {$G_{n\bar n}$};
\end{tikzpicture}
\caption{SCET 的模式（\cref{def:scet-modes}）在 $(\bar n\cdot p,\,n\cdot p)$ 对数平面中的位置，坐标为 $1/\lambda$ 的幂次（示意）。在壳模式位于等虚度线 $p^2\simeq\bar n\cdot p\,n\cdot p$ 上（虚线，自左下至右上依次为 $p^2\sim\lambda Q^2$、$\lambda^2Q^2$、$\lambda^4Q^2$）；点线是快度 $y=\frac12\ln(\bar n\cdot p/n\cdot p)=0$ 的点，垂直于它的方向改变快度，$n$ 共线模式的快度为正。空心记号为 Glauber 模式 $G_{n\bar n}$、$G_{ns}$、$G_{\bar ns}$：光锥分量与所在位置相同，但 $p_\perp\sim\lambda Q$，故 $p^2\simeq-\bm p_{\mathrm T}^2$ 远离质壳。SCET$_{\rm I}$：硬、两个共线、超软；SCET$_{\rm II}$：硬、两个共线、软。}
\label{fig:scet-modes}
\end{figure}
%%% END fig:scet-modes

\begin{proposition}[场的标度 \tagR]\label{prop:scet-scalings}
（$R_{\mathrm{def}}$：由算符定义直接推出。）以 $Q$ 的相应幂次为单位（夸克场 $Q^{3/2}$，胶子场 $Q$），各模式的场按下列方式标度：
\begin{equation}\label{eq:scet-fieldscaling}
\xi_n\sim\lambda,\qquad (\bar n\cdot A_n,\,n\cdot A_n,\,A_{n\perp})\sim(1,\lambda^2,\lambda),\qquad
A_{us}\sim\lambda^2,\quad q_{us}\sim\lambda^3,\qquad A_s\sim\lambda,\quad q_s\sim\lambda^{3/2}.
\end{equation}
% VF-SCET-2
（本命题只涉及部分子场，不涉及强子态。）
\uses{def:scet-modes, def:prelim-feynman-rules}
\end{proposition}

\begin{proof}
场 $\phi$ 的标度定义为使自由两点函数 $\langle0|\mathrm T\phi(x)\phi(0)|0\rangle=\int\frac{\dd^4p}{(2\pi)^4}\ee^{-\ii p\cdot x}\Delta_\phi(p)$ 在该模式的动量区内与 $\phi\phi$ 同阶的幂次，即 $\phi\phi\sim[\dd^4p]\,\Delta_\phi$，其中 $\dd^4p=\frac12\dd(\bar n\cdot p)\,\dd(n\cdot p)\,\dd^2p_\perp$。共线模式 $[\dd^4p]\sim1\cdot\lambda^2\cdot\lambda^2=\lambda^4$；共线夸克场 $\xi_n=P_n\psi$ 的传播子是 QCD 传播子的投影：把 $\sla p=\frac{\bar n\cdot p}{2}\sla n+\frac{n\cdot p}{2}\bar{\sla n}+\sla p_\perp$ 代入并用 $P_n\sla nP_{\bar n}=\sla n$、$P_n\bar{\sla n}P_{\bar n}=0$、$P_n\gamma_\perp^\mu P_{\bar n}=0$，得
\begin{equation}\label{eq:scet-xiprop}
P_n\,\frac{\ii\sla p}{p^2+\ii0}\,P_{\bar n}=\ii\,\frac{\sla n}{2}\,\frac{\bar n\cdot p}{p^2+\ii0}\sim\lambda^{-2}
\end{equation}
（对离壳的 $p$ 也精确成立），故 $\xi_n\xi_n\sim\lambda^2$，$\xi_n\sim\lambda$。共线胶子在一般协变规范中的传播子 $\frac{-\ii}{k^2}[g^{\mu\nu}-(1-\xi_g)\frac{k^\mu k^\nu}{k^2}]$（\cref{def:prelim-feynman-rules}）的第二项给出 $A_n^\mu A_n^\nu\sim\lambda^4\cdot k^\mu k^\nu/\lambda^4=k^\mu k^\nu$，即 $A_n^\mu$ 的各分量与 $k^\mu$ 的相应分量同阶；第一项给出 $\langle\bar n\cdot A_n\,\bar n\cdot A_n\rangle\propto\bar n^2=0$、$\langle n\cdot A_n\,\bar n\cdot A_n\rangle\propto n\cdot\bar n/k^2\cdot\lambda^4\sim\lambda^2$、$\langle A_{n\perp}A_{n\perp}\rangle\sim\lambda^2$，与 $(1,\lambda^2,\lambda)$ 相容。超软模式 $[\dd^4k]\sim\lambda^8$，胶子与夸克的传播子分别 $\sim\lambda^{-4}$ 与 $\sla k/k^2\sim\lambda^{-2}$，故 $A_{us}\sim\lambda^2$，$q_{us}\sim\lambda^3$；软模式 $[\dd^4k]\sim\lambda^4$，传播子 $\sim\lambda^{-2}$ 与 $\lambda^{-1}$，故 $A_s\sim\lambda$，$q_s\sim\lambda^{3/2}$。
\end{proof}

于是协变导数 $\ii D^\mu=\ii\partial^\mu+gA^\mu$ 的每个分量都与相应的动量分量同阶；特别地 $\bar n\cdot A_n\sim\lambda^0$ 与 $\bar n\cdot p$ 同阶，这正是 $g\,\bar n\cdot A_n$ 必须对所有阶求和（共线 Wilson 线，\cref{prop:scet-wilson}）的原因。

\paragraph{模式之间的相互作用与多极展开。}
设 $p$ 为在壳的 $n$ 共线动量，由 $2p\cdot k=\bar n\cdot p\,n\cdot k+n\cdot p\,\bar n\cdot k+2p_\perp\cdot k_\perp$ 对各模式逐项计数（全部组合我们都逐一核对过）：% w18_s1 (E)
(i) 吸收超软动量后 $(p+k)^2=p^2+\bar n\cdot p\,n\cdot k+\order{\lambda^3Q^2}\sim\lambda^2Q^2$，仍是共线的，且在领头幂次只出现 $n\cdot k$；(ii) 吸收软动量后 $(p+k)^2\simeq\bar n\cdot p\,n\cdot k\sim\lambda Q^2$，是硬共线虚度，这一中间线须积出（\cref{sec:scet-scet2}）；(iii) 与 $\bar n$ 共线动量结合得 $\sim Q^2$，是硬的，积出后给出匹配系数与 $W_n$（\cref{prop:scet-wilson}）；(iv) 交换 Glauber 动量 $q$ 时两个共线扇区都保持在壳，而 $|n\cdot q\,\bar n\cdot q|\ll\bm q_{\mathrm T}^2$，Glauber 交换在 $x^\pm$ 上是瞬时的“势”（\cref{sec:scet-glauber}）。共线场所在的时空区域满足 $(n\cdot x,\bar n\cdot x,x_\perp)\sim(1,\lambda^{-2},\lambda^{-1})/Q$，因此
\begin{equation}\label{eq:scet-multipole}
k_{us}\cdot x=\tfrac12\,n\cdot k\,\bar n\cdot x\,[\lambda^0]+k_\perp\cdot x_\perp\,[\lambda]+\tfrac12\,\bar n\cdot k\,n\cdot x\,[\lambda^2],\qquad
A_{us}(x)=A_{us}(x_-)+\order{\lambda},\quad x_-^\mu\equiv\bar n\cdot x\,\frac{n^\mu}{2}:
\end{equation}
在领头幂次，共线场看到的超软场只依赖于 $\bar n\cdot x$（多极展开）。等价地，可以把共线动量分为大“标签” $(\bar n\cdot p,p_\perp)$ 与 $\sim\lambda^2Q$ 的剩余动量，用标签算符 $\bar{\mathcal P}$、$\mathcal P_\perp$ 表述\cite{BauerEtAl2001,BauerStewart2001}；坐标空间的等价表述见\cite{BenekeEtAl2002,Beneke:2002ni}。下面采用坐标空间的写法：$\ii\bar n\cdot\partial$ 与 $\ii\partial_\perp$ 作用于共线场时只取出大分量。领头阶共线拉氏量的每一项都是 $\order{\lambda^4}$，而共线区域的测度 $\dd^4x\sim\lambda^{-4}$，所以共线作用量是 $\order{\lambda^0}$；超软部分相应为 $\lambda^8$ 与 $\lambda^{-8}$。
